\clearpage
\fancyhead[L]{\small DETERMINING ABSOLUTE DECLINATION.}
\SourceLeaf{442}\pdfbookmark[0]{On using the magnetometer to determine absolute declination.}{absolute-declination-title}
\begin{center}
{\Large\bfseries ON USING THE MAGNETOMETER\par\medskip TO DETERMINE ABSOLUTE DECLINATION.}
\par\bigskip
Resultate aus den Beobachtungen des magnetischen Vereins. 1841. I.
\end{center}
I do not intend to treat the subject named in the title in full once again. A very detailed paper on it has already appeared in volume 2 of the \textit{Resultate} [1837. VII.]. I shall limit myself here to a single main problem for which the treatment given there on pp. 121—124 seems inadequate. The problem is to determine the azimuth of the vertical plane containing the optical axis of the observing telescope.

The vertical plane in question is fixed by the mark and a fixed point on the scale, which is determined by the plumb line hanging above the centre of the objective of the observing telescope. A distant object whose azimuth is already known from other sources must be visible from the position of the observing telescope. The first task is therefore to determine the angle between this object and the mark, projected onto the horizontal plane. I assume that a theodolite of the familiar design introduced by \textsc{Reichenbach} is used for this task. This leaves open whether the same theodolite is used for the magnetic observations, for which a separate reading telescope can readily be used.

Such theodolites are ordinarily used to measure an\-gles between objects at such great distances that a slight
\expandafter\def\csname HistoricalNoteText443\endcsname{About 20 millimetres in the observations to be given below.}%
\SourceLeaf{443}departure in the actual construction from several of the conditions on which the design is based cannot produce an appreciable error. Indeed, absolute perfection is unattainable in any mechanical work. But when the objects (or one of them, as in the present case) are comparatively very close, these departures do require closer attention, and the following circumstances in particular must be considered.

I. The vertical axis of rotation, the horizontal axis of rotation and the optical axis of the telescope should intersect at a single point. Insofar as this condition is imperfectly met, there will be a threefold departure. Let $A,B$ be the two points on the vertical and horizontal axes of rotation, respectively, where these axes are closest to each other. Similarly, let $C,D$ be the corresponding points on the horizontal axis of rotation and the optical axis. Denote the distances $AB,BC,CD$ by $\alpha,\beta,\gamma$, with any convention for their signs. For example, one may take $\alpha$ as positive if $A$ is on the same side of the horizontal axis of rotation as the telescope’s eyepiece, $\beta$ as positive if $C$ lies to the right of $B$ for the observer at the eyepiece, and $\gamma$ as positive if $D$ lies above $C$.

II. The optical axis of the telescope should be perpendicular to the horizontal axis of rotation. This condition can be satisfied with all the required precision. However, to see a nearby object clearly after observing a distant one, the eyepiece tube must be drawn farther\HistoricalNote{443} out, giving the crosshairs a different position relative to the objective. We therefore have no grounds for assuming that the two positions of the eyepiece tube correspond to the same optical axis and must allow for the adjustment made for one position to be lost in the other. For greater generality, we may suppose that the adjustment has been made exactly for neither position and denote the collimation error by $c$ for the first position and by $c'$ for the second. These errors may be taken as positive when the optical axis makes an acute angle with the arm of the horizontal axis that lies to the observer’s right.

\SourceLeaf{444}Clearly, if the quantities $\beta,\gamma$ belong to the first eyepiece position, somewhat different values will take their place in the second position. We may denote these by $\beta',\gamma'$.

It is easy to express the effect of all these deviations on the measurement of both the horizontal angle between the two objects and their angles of elevation (if the theodolite also has a vertical circle) in exact equations, from which the results would be obtained by means of an equation of the fourth degree. But since all seven quantities $\alpha,\beta,\gamma,\beta',\gamma',c,c'$ can only be very small, we can safely neglect all quantities. that are of second or higher order in relation to them and express the result of their influence in a very simple form. We can dispense even with this calculation here. It is easy to see that if the telescope is reversed in the proper way, all seven deviations simply take the opposite signs while retaining their magnitudes. The same will therefore be true of the errors in the measurements made with the telescope seated in the two different ways. The mean of these two measurements is thus automatically free of the influence of these errors, without any need to know their separate components. We thereby obtain the true value of the angle between the two vertical planes that intersect in the vertical axis of the instrument and contain the two objects. The same is true of the elevations, which then refer to the point $A$, but are unnecessary for our present purpose.

The reversal must be carried out so that the pivots return to the same bearings, while the upper side of the telescope becomes the lower side and the objective takes the place of the eyepiece. This reversal thus amounts to a half-turn about the horizontal axis, for which the supports are simply too low. If instead the telescope were reversed by placing the pivots in the opposite bearings while keeping the objective end on the same side (which geometrically amounts to a half-turn about the axis of the telescope), the seven quantities $\alpha,\beta,\gamma,\beta',\gamma',c,c'$ would no longer all simply take the opposite signs. Only $\gamma,\gamma',c,c'$ would simply change sign. There is no assurance that the supports are \textit{exactly}
\expandafter\def\csname HistoricalNoteText445\endcsname{This tower is indicated on the first plate of the first volume of the \textit{Resultate} [1836. I.], roughly as it appears from the theodolite station with the window closed. From the position of the eye on which the perspective drawing is actually based, the tower is hidden by the wall to the left of the window.}%
\SourceLeaf{445}equidistant from the vertical axis, and hence. after such a reversal, the point $B$ might be different from the one before, so $\beta$ and $\beta'$ could also take different values. The fact that $\alpha$ keeps exactly its former value instead of changing sign is immaterial here, since $\alpha$ does not occur at all in the linear expression for the error in the measurement of the horizontal angle.

The simplest way to show how such an angle measurement can be used for the intended purpose is by an example. For this I choose the most recent application of the procedure, carried out on 11 March of this year.

At the local magnetic observatory, a church tower in the town is used to tie the observations to the true astronomical meridian. The tower’s finial rod is clearly visible from the station of the observing telescope through the open north window\HistoricalNote{445}, from the centre of the pier that has stood since July 1837 in place of the wooden tripod formerly used, at an azimuth of $173^0 35'25''5$. This azimuth was found by setting up a theodolite at another point in the room, with its vertical axis exactly in line with the centre of the pier and the church tower, and measuring the angles between that tower and two other church towers visible from there. The positions of these towers relative to the zero point at the astronomical observatory were known precisely from earlier measurements tied to the arc survey, so the azimuth in question could easily be calculated from those angle measurements.

An eight-inch \textsc{Ertel} repeating theodolite was then set up on the pier with its vertical axis passing as nearly as possible through the centre of the pier. The horizontal angle between the mark and the tower’s finial rod was measured with the telescope seated in the two ways described above, by 25 repetitions in each case. In the first position the angle was found to be
\[
=11^0 40'54''50
\]
\SourceLeaf{446}in the second
\[
=11^0 41'36''18
\]

The true value of the angle, with its vertex placed in the vertical axis of the theodolite, is therefore
\[
=11^0 41'15''34
\]
and the azimuth of the vertical plane through this vertical axis and the mark is consequently
\[
=161^0 54'10''16
\]

Another operation was carried out in conjunction with this one to find the point where this vertical plane cuts the scale.

A ring is fitted over the objective end of the theodolite telescope. On its front face it has two fine, diametrically opposite notches, with two corresponding small hooks on the rounded outer surface. A fine gold thread with a weight attached is hung from one or the other hook, depending on the position of the telescope. The ring is turned so that the diameter through the notches is perpendicular to the horizontal axis of rotation of the telescope. We can tell that this condition is met when the plumb line lying in the upper notch also aligns exactly with the lower one. To check this, the telescope must be set nearly horizontal, inclined downwards just enough for the thread still to hang freely in front of the ring. The coincidence is checked with a magnifying glass. The plumb line hangs a very short distance in front of the scale. The task is now to record the corresponding points on the scale in the two different positions of the telescope, with the telescope directed each time at the mark, or more precisely in its vertical plane. Strictly speaking, these are the points on the scale that lie in the vertical plane through the mark and the plumb line. They can be seen directly in the magnetometer mirror when the theodolite itself is to serve as the reading telescope, so that the scale is at a height suited to that use. It is hardly necessary to point out that in this case it may be necessary first to bring the bar magnet
\SourceLeaf{447}of the magnetometer into a position in which the relevant scale point appears near the crosshairs of the theodolite telescope, by means of a deflecting bar acting from a distance. At the local magnetic observatory, magnetic observations are now made with a separate reading telescope, which is lower above the pier than the theodolite telescope. Through the latter, the image of the scale in the magnetometer mirror is not visible because the scale is set at a height suited to the position of the reading telescope. To determine the scale point corresponding to the plumb line hanging from the theodolite telescope, I therefore used the reading telescope itself, set up for this purpose close to the mark in the vertical plane in question. There is clearly no need to be overly scrupulous about this last condition, since the plumb line, as already mentioned, hangs only a short distance from the scale. The plumb line was thus found to correspond
\begin{quote}
to scale point $850.8$ in the theodolite telescope’s first position, and

 to point $849.4$ in the second position,
\end{quote}
from which we may conclude that the vertical plane passing through the mark and the vertical axis of the theodolite, whose azimuth was determined above, cuts the scale at point $850.1$.

Determining the azimuth of the vertical plane containing the optical axis of the observing telescope now presents no further difficulty. If the plumb line hanging in front of the centre of its objective corresponds to scale point $850{,}1+n$, it suffices (since the scale is assumed to be set perpendicular to that plane) to divide the product $n.206265''$ by the horizontal distance of the scale from the mark, expressed in scale divisions, and to add the quotient with its sign to $161^0 54'10''16$. At present that distance is $=9638.7$. If the theodolite itself were used for observing, in the first position of the telescope this azimuth would be
\[
=161^0 54'25''1
\]
and in the second position
\[
=161^0 53'55''2
\]
\SourceLeaf{448}But, as already noted, a separate reading telescope is used for observing. After the preceding operations it was set up so that, with the optical axis directed at the vertical through the mark, the plumb line corresponded to point $850.0$. The required azimuth is therefore
\[
=161^0 54'8''0
\]

A few remarks about the procedure discussed here may be added.

I. If the horizontal axis has some play along its length in its bearings, care must be taken that it always has the same position relative to the supports in the individual angle measurements. One way is to take up the play on one side each time by pressing lightly against the end of a particular pivot. Without this precaution, one could not count on the quantity denoted above by $\beta$ retaining the same value in all repetitions in the first position of the telescope and always having exactly the opposite value in the second.

II. The operations described here do not require the optical axis of the theodolite telescope to be adjusted precisely for one of the two eyepiece positions, that is, to be perpendicular to the horizontal axis of rotation. If the theodolite also serves as a reading telescope, however, this adjustment must certainly be made before it is used in that way. It must be made for the position of the eyepiece tube used for observing, when the mark and the reflected image of the scale appear clearly. As is well known, the perpendicularity of the optical axis to the horizontal axis of rotation is checked by reversing the telescope, and specifically by the reversal that could \textit{not} be used in the angle measurements above [p. 439 of this volume], namely by putting the pivots in the opposite bearings without changing the direction in which the telescope points. Such a check usually concerns the eyepiece-tube position in which very distant objects can be seen clearly. and in this case it is indeed enough for such an object to appear on the crosshairs both before and after reversal. In the present case, however, if the plumb line hanging in front of the centre of the objective is in a different position after reversal, a second aiming point must be used next to the first, displaced by the same amount. Clearly, a separate reading telescope used in place of the theodo%
\SourceLeaf{449}%
lite must undergo the same adjustment and must therefore be mounted suitably for it. The horizontal axis of rotation must, of course, also be properly levelled. The two reading telescopes used with the local magnetometers, though of considerably greater optical power than is usual for theodolite telescopes, are mounted in almost exactly the same way as theodolites, only without graduated circles.

It may also be noted that the \textit{effect} of a collimation error on the azimuth of the optical axis is only a very small fraction of the collimation error itself. That fraction is determined by the difference between the secants of the two inclinations when the telescope is directed first at the mark and then at the mirror. In the local unifilar magnetometer these inclinations are $1^0 55'$ and $5^0 16'$. The difference between the azimuths of the optical axis when directed at the mark and at the mirror is therefore only $\frac{1}{272}$ of the collimation error itself. In similar circumstances it will therefore usually suffice to adjust the collimation using a distant object. For unless this adjustment has moved the crosshairs far from the centre of the eyepiece tube, drawing the tube farther out can hardly produce a collimation error of more than a small fraction of a minute of arc, so its effect remains quite imperceptible.

III. The purpose of leaving the plumb line permanently hanging from the observing telescope is to make any accidental displacement of the scale immediately apparent. If this occurs, one may either return the scale to its former position or count in the calculation from the point on the scale that corresponds to the plumb line after the change, in the same way as one previously counted from the former point. With the way the scale is now fastened to the pier at the magnetic observatory, however, accidental displacements no longer occur at all.

